\begin{exsol}{21}{Tadelis Exercise 9.4}{188}{Exam}{L4}
\asked A Cournot stage ($p = 100 - q$, zero cost) is followed by an ``access'' game with
equilibria $(O,o)$ and $(N,n)$. (a) The Cournot equilibrium, the access game's pure equilibria and
all pure SPE when $\delta = 0$; (b) an SPE with monopoly output split equally when $\delta = 1$;
(c) the lowest $\delta$ for it; (d) whether it survives at $\delta = 0.4$, and if not the best
symmetric profits.

\method Multistage games with a carrot and a stick in the last stage (Tadelis \S9.4,
pp.~180--184; threshold $\delta^* = g/(c - s)$, \texttt{Notas/L04/L04-02}); Cournot best replies
(Tadelis \S5.2.3, pp.~87--88). Payoffs $v_i = \pi_i(q_1, q_2) + \delta\,u_i(\text{access})$.

\solution
\textbf{(a) Stage games and $\delta = 0$.}

\emph{Cournot.} Firm $i$ maximizes $(100 - q_i - q_j)\,q_i$:
\begin{equation}
  100 - 2q_i - q_j = 0 \;\Rightarrow\; q_i = BR_i(q_j) = \frac{100 - q_j}{2}. \label{eq:21br}
\end{equation}
By symmetry $q_i = q_j = q$, so $q = \frac{100 - q}{2}$, i.e.\ $q^C = \frac{100}{3}$, price
$\frac{100}{3}$, and $\pi^C = \bigl(\frac{100}{3}\bigr)^2 = \frac{10000}{9} \approx 1111.1$ each.

\emph{Access game.} Payoffs (player~1 rows):
\begin{center}
\begin{tabular}{r|c|c|}
\multicolumn{1}{r}{} & \multicolumn{1}{c}{$o$} & \multicolumn{1}{c}{$n$}\\ \cline{2-3}
$O$ & $2500,\ 2500$ & $-10,\ 0$ \\ \cline{2-3}
$N$ & $0,\ -10$ & $0,\ 0$ \\ \cline{2-3}
\end{tabular}
\end{center}
$O$ is the best reply to $o$ ($2500 > 0$) and $N$ the best reply to $n$ ($0 > -10$); symmetrically
for player~2. Pure equilibria: $(O,o)$ with $(2500, 2500)$ and $(N,n)$ with $(0,0)$.

\emph{SPE with $\delta = 0$.} Stage-2 payoffs carry weight $0$, so continuation play cannot reward or
punish anything, and stage~1 must be the Cournot equilibrium. The pure SPE are
\[
  s_i:\ q_i = \tfrac{100}{3};\ \text{in stage 2, after each } (q_1, q_2),\ \text{play the action
  of } e(q_1, q_2),
\]
where $e$ is \emph{any} map from $\mathbb{R}_+^2$ to $\{(O,o), (N,n)\}$ (e.g.\ always $(O,o)$).
There is a continuum of them, all with first-stage outcome $(\frac{100}{3}, \frac{100}{3})$.
Strictly, with $\delta = 0$ the stage-2 actions do not affect any payoff, so every stage-2 profile
is an equilibrium of those subgames; the answer above uses the standard reading that stage-2
actions are judged by stage-2 payoffs (the limit $\delta \to 0^+$).

\textbf{(b) Monopoly split with $\delta = 1$.} Joint profit $(100 - Q)Q$ is maximized at $Q = 50$;
each firm produces $25$ and earns
\[
  \pi^M = 25\,(100 - 50) = 1250 .
\]
Strategies (for both firms):
\[
  s_i:\ q_i = 25;\ \text{in stage 2, open } (O / o) \text{ if } (q_1,q_2) = (25,25),\ \text{keep
  closed } (N / n) \text{ otherwise}.
\]
\emph{Stage 2.} After every history the prescribed profile is an equilibrium of the access game.
\emph{Stage 1.} The best deviation against $q_j = 25$ is, by \eqref{eq:21br},
$q_i = \frac{100 - 25}{2} = 37.5$, with price $100 - 62.5 = 37.5$ and
\[
  \pi^D = 37.5 \times 37.5 = 1406.25 .
\]
Staying gives $1250 + \delta\,(2500)$; deviating gives $1406.25 + \delta\,(0)$. With $\delta = 1$:
$3750 \ge 1406.25$. It is an SPE.

\textbf{(c) Lowest $\delta$.}
\begin{equation}
  1250 + 2500\,\delta \;\ge\; 1406.25 \iff \delta \;\ge\; \delta^* = \frac{156.25}{2500} = \frac{1}{16} = 0.0625 .
  \label{eq:21d}
\end{equation}
In the carrot-and-stick form: $g = 156.25$, $c - s = 2500 - 0$.

\textbf{(d) $\delta = 0.4$.} Since $0.4 > \frac1{16}$, condition \eqref{eq:21d} holds:
$1250 + 0.4\,(2500) = 2250 \ge 1406.25$. The monopoly split \emph{is} supported, with the same
strategies as in (b), and each firm earns $2250$ in discounted total. The ``if not'' branch does not
arise with these numbers.

For completeness, when $\delta < \frac1{16}$: with symmetric output $q$ and the same threat, the
deviation gain is
\[
  g(q) = \Bigl(\frac{100 - q}{2}\Bigr)^2 - q\,(100 - 2q) = \frac94\Bigl(\frac{100}{3} - q\Bigr)^2 ,
\]
and $g(q) \le 2500\,\delta$ holds iff $q \ge \frac{100}{3}(1 - \sqrt{\delta})$. The best symmetric
equilibrium uses $q = \max\{25,\ \frac{100}{3}(1 - \sqrt\delta)\}$; e.g.\ at $\delta = 0.04$,
$q = \frac{80}{3}$ and stage-1 profit $\frac{11200}{9} \approx 1244.4$ each.

\interp The second stage is worth far more ($2500$) than the one-shot gain from cheating on the
cartel ($156.25$), so even very impatient firms ($\delta \ge 0.0625$) can collude in stage~1. The
threat is credible because $(N,n)$ is itself an equilibrium: once someone cheats, neither firm
wants to be the only one to open. A future coordination game with a good and a bad equilibrium is
an effective enforcement device, which is the multistage logic of Tadelis \S9.4.

\verify \texttt{check\_B\_L04.py} recomputes the Cournot equilibrium ($\frac{100}{3}$,
$\frac{10000}{9}$), the two pure equilibria of the access game, $\pi^M = 1250$,
$\pi^D = 1406.25$, $\delta^* = \frac1{16}$, the total $2250$ at $\delta = 0.4$, the identity for
$g(q)$, the bound $q \ge \frac{100}{3}(1 - \sqrt\delta)$ (equal to $25$ at $\delta = \frac1{16}$),
and the $\delta = 0.04$ example. Tadelis prints no answer.

\answer{(a) Cournot $q_i = \frac{100}{3}$, $\pi_i = \frac{10000}{9}$; access game pure equilibria
$(O,o)$ and $(N,n)$; with $\delta = 0$ every pure SPE has $q_i = \frac{100}{3}$ followed by
$(O,o)$ or $(N,n)$ after each history (any assignment). (b) $q_i = 25$; open after $(25,25)$, closed
otherwise. (c) $\delta^* = \frac1{16}$. (d) Yes, $0.4 > \frac1{16}$: monopoly split supported,
discounted payoff $2250$ each.}
\end{exsol}
